\chapter{Dati e pipeline}
\label{cap:2}

Questo capitolo descrive come il dato arriva dalle fonti, come si controlla e come diventa confrontabile fra domini.

\section{Fonti}
\label{sec:dati-fonti}

La Tabella~\ref{tab:fonti} elenca le fonti usate. Sono tutte pubbliche e gratuite: OpenAQ e i servizi Copernicus richiedono solo una chiave gratuita.

{\footnotesize
\setlength{\tabcolsep}{3pt}
\begin{longtable}{>{\raggedright\arraybackslash}p{1.5cm}>{\raggedright\arraybackslash}p{2.7cm}>{\raggedright\arraybackslash}p{2.9cm}>{\raggedright\arraybackslash}p{2.2cm}>{\raggedright\arraybackslash}p{1.5cm}>{\raggedright\arraybackslash}p{1.6cm}}
\caption{Fonti dei dati usate dal progetto.}
\label{tab:fonti}\\
\hline
\textbf{Dominio} & \textbf{Fonte} & \textbf{Dato} & \textbf{Grana} & \textbf{Periodo} & \textbf{Licenza}\\
\hline
\endfirsthead
\hline
\endfoot
Aria & ARPA Puglia~\cite{arpa_puglia_api} & misure della rete regionale & stazione $\times$ giorno & 2020--oggi & CC BY 4.0\\
 & EEA~\cite{eea_aqds} & misure delle stesse stazioni & stazione $\times$ ora & 2020--oggi & CC BY 4.0\\
 & OpenAQ~\cite{openaq} & misure riprese dall'EEA & stazione $\times$ ora & 2020--oggi & ODC-BY\\
 & Copernicus CAMS~\cite{cams_reanalisi,cams_previsioni} & rianalisi e analisi europea (modellato) & cella 0,1° $\times$ ora & 2020--oggi & CC BY 4.0\\
\hline
Calore & ERA5-Land~\cite{era5_land} & temperatura media, massima, minima (modellato) & cella 0,1° $\times$ giorno & 2024--oggi & CC BY 4.0\\
 & ARPA Puglia meteo~\cite{arpa_meteo_rrqa} & temperatura misurata, 15 stazioni & stazione $\times$ ora & 2020--oggi & CC BY 4.0\\
 & CLMS~\cite{copernicus_tcd,copernicus_imperviousness} & copertura arborea (2024), suolo impermeabile (2018) & raster 10\,m & statico & libero, con citazione\\
\hline
Salute & ISTAT~\cite{istat_dati_mortalita} & decessi per tutte le cause & comune $\times$ giorno & 2011--2026 & CC BY 4.0\\
 & ISTAT~\cite{istat_bilancio_mensile} & bilancio demografico (validazione) & comune $\times$ mese & 2019--2024 & CC BY 4.0\\
 & Eurostat~\cite{eurostat_decessi_settimanali} & decessi (validazione) & provincia $\times$ settimana & 2011--2026 & CC BY 4.0\\
 & ISTAT Health for All~\cite{istat_hfa} & dimissioni ospedaliere (informative) & provincia $\times$ anno & 1999--2024 & CC BY 4.0\\
\end{longtable}
}

\section{Acquisizione e qualità del dato}
\label{sec:dati-qualita}

I dati si prelevano con chiamate dirette alle API delle fonti, senza database. L'archivio è incrementale, con un file per mese: i mesi chiusi non si riscaricano, il mese in corso si rigenera ogni giorno. Ogni fonte è un adattatore con la stessa interfaccia. Un errore di una fonte non interrompe la generazione: le celle coinvolte restano vuote e l'errore va in un registro.

Ogni cella contiene un valore, uno stato e una nota. Lo stato è \texttt{VALORE} se il dato è valido, \texttt{NULL} se la fonte copre la cella ma il valore non è utilizzabile, \texttt{NO\_DATA} se la fonte non la copre (per esempio un comune senza stazioni). I valori si controllano appena ricevuti: campi attesi, numero finito e non negativo, nessuna riga discordante per la stessa cella. Il motivo di un vuoto resta visibile anche dopo l'aggregazione. Un valore implausibile (oltre la soglia di allarme di legge o dieci volte il valore limite) esce dalla media per comune ma resta nel dataset.

Il valore giornaliero di EEA, OpenAQ e CAMS è la media delle ore del giorno, con la completezza indicata nella nota. ARPA pubblica invece un valore già giornaliero, che per NO\textsubscript{2}, O\textsubscript{3}, SO\textsubscript{2} e CO è il massimo: per questo resta una serie a sé.

\section{Validazione per ridondanza}
\label{sec:dati-validazione}

Fonti indipendenti che descrivono la stessa grandezza si confrontano sulle stesse celle (Tabella~\ref{tab:validazione}).

{\footnotesize
\setlength{\tabcolsep}{3pt}
\begin{longtable}{>{\raggedright\arraybackslash}p{3.0cm}>{\raggedright\arraybackslash}p{2.7cm}>{\raggedright\arraybackslash}p{7.6cm}}
\caption{Validazione per ridondanza: risultati principali.}
\label{tab:validazione}\\
\hline
\textbf{Confronto} & \textbf{Grana} & \textbf{Risultato}\\
\hline
\endfirsthead
\hline
\endfoot
ARPA -- EEA & stazione $\times$ giorno, giugno 2024 & PM\textsubscript{10}: 100\% delle celle entro 1\,µg/m\textsuperscript{3}. Per gli inquinanti di cui ARPA pubblica il massimo il confronto è solo qualitativo.\\
EEA -- CAMS & comune $\times$ giorno, giugno 2024 & correlazione 0,93 per PM\textsubscript{10}, 0,90 per PM\textsubscript{2.5}, 0,83 per O\textsubscript{3}; solo 0,31 per NO\textsubscript{2}, che il modello stima più basso di circa 4\,µg/m\textsuperscript{3}.\\
ERA5-Land -- ARPA meteo & stazione $\times$ giorno, 2024--2026 & correlazione 0,97--0,99. Media ben riprodotta (scarto $-$0,15\,°C); massima sottostimata ($-$1,1\,°C) e minima sovrastimata (+0,5\,°C): meno giorni oltre i 35\,°C e più notti tropicali.\\
ISTAT -- Eurostat & provincia $\times$ settimana & identici al 100\% dal 2011 al 2023.\\
ISTAT -- bilancio mensile & comune $\times$ mese, 2019--2024 & identici nel 74--97\% delle celle, entro $\pm$1 decesso nel 94--99,8\%.\\
\end{longtable}
}
\paragraph{Limiti noti.} Alcuni casi non si correggono e si dichiarano.
\begin{itemize}
  \item \textbf{Fonti non uniformi}: due stazioni di OpenAQ (provider HabitatMap, solo PM\textsubscript{2.5}) sono sensori di cittadini, senza licenza dichiarata dall'API. Sono incluse e segnalate, e non equivalgono alle altre, che sono ODC-BY.
  \item \textbf{Valori anomali}: la soglia di plausibilità esclude dalla media del comune i valori impossibili, ma non quelli solo insoliti. Un SO\textsubscript{2} di 65,6\,µg/m\textsuperscript{3} alla stazione di Torchiarolo resta nel dato.
  \item \textbf{Orari e unità}: OpenAQ pubblica alcune serie con orari sfasati, e si converte l'orario dichiarato senza correggerlo. Il CO cambia unità a maggio 2026 e si lascia com'è.
  \item \textbf{Copertura temporale}: i decessi ISTAT arrivano al 30 giugno 2026, perché è il limite della fonte.
\end{itemize}

\section{Armonizzazione spaziale}
\label{sec:dati-armonizzazione}

\begin{itemize}
  \item \textbf{Stazioni}: quelle di EEA e OpenAQ si abbinano alle stazioni ARPA per coordinate (entro 250\,m, o 1\,km se il comune compare nel nome). Ogni stazione va al comune dichiarato dalla fonte, altrimenti a quello che la contiene.
  \item \textbf{Griglie} (CAMS, ERA5-Land): il valore del comune è la media delle celle, pesata sulla quota di superficie comunale in ciascuna. CAMS diventa una stazione virtuale per comune, esclusa dalla media delle stazioni fisiche.
  \item \textbf{Raster} (CLMS): media dei pixel dentro il confine comunale.
  \item \textbf{Vista «Tutte»}: per comune e giorno si mediano le sole fonti concordi, entro il 20\% (o 3\,µg/m\textsuperscript{3}) dalla mediana; ARPA entra solo per il particolato. Dove mancano le stazioni resta CAMS. Il dataset resta separato per fonte.
\end{itemize}

\section{Generazione dei dati simulati}
\label{sec:dati-simulazione}

I dati sanitari simulati (Sezione~\ref{sec:domini-salute}) si generano per comune e per giorno, poi si sommano al distretto.
\begin{itemize}
  \item \textbf{Livello annuo}: ricoveri dalle dimissioni Health for All per provincia, ripartite sui comuni per residenti; pronto soccorso dal tasso nazionale di 311 accessi ogni 1.000 abitanti~\cite{qs_annuario_ssn_2023}; ADI da una coorte sintetica con fragilità individuale, per cui pochi assistiti tornano spesso al pronto soccorso~\cite{campagna_adi_ps_2021}.
  \item \textbf{Stagionalità}: profilo medio dei decessi ISTAT reali del Salento 2011--2019, più un effetto ipotetico del giorno della settimana.
  \item \textbf{Effetto dell'ambiente} (solo scenario \emph{letteratura}): +0,75\% per 10\,µg/m\textsuperscript{3} di PM\textsubscript{10} sugli esiti respiratori, +0,57\% per 10\,µg/m\textsuperscript{3} di NO\textsubscript{2} su quelli cardiovascolari~\cite{scarinzi_epiair2_2013}, +4,5\% per grado di temperatura massima oltre i 30\,°C sugli esiti respiratori~\cite{michelozzi_phewe_2009}. L'esposizione viene da CAMS ed ERA5-Land, che coprono tutti i comuni.
  \item \textbf{Conteggi}: distribuzione di Poisson con seme fisso, lo stesso nei due scenari, così la differenza dipende solo dall'effetto.
\end{itemize}

Nella verifica i ricoveri simulati restano entro l'1,5\% dalle dimissioni Health for All 2024 per provincia. Un modello di Poisson standard ritrova l'effetto del caldo imposto (stime fra +4,7\% e +5,3\%) e nessun effetto nello scenario nullo. Gli effetti di PM\textsubscript{10} e NO\textsubscript{2}, inferiori all'1\%, non si distinguono dal caso: è un limite di potenza statistica che varrà anche per i dati reali.
